\section{Two-torsion excludes the exponent prime from the support}
\label{ts-section}

The fixed boundary curve has the nonzero rational two-torsion point
$(0,0)$. At every split rational prime $q>3$, good reduction preserves
it, so
\begin{equation}\label{ts-even-trace}
 t_q(E_0)=q+1-\#E_0(\mathbb F_q)\quad\text{is even}.
\end{equation}
This elementary parity condition gives a stronger consequence of the
boundary module comparison, including at the exponent prime itself.

\begin{theorem}[Boundary support with no exponent-prime exception]
\label{ts-boundary-support}
Let $a,b$ be positive coprime integers and $p>7$ prime. Let $\nu$
be a quadratic character of $G_K$ unramified outside the primes over
$6$. Suppose
\begin{equation}\label{ts-module}
 E_{a,b}[p]\otimes\overline{\mathbb F}_p
 \simeq(E_0[p]\otimes\overline{\mathbb F}_p)\otimes\nu.
\end{equation}
Then $p\nmid F(a,b)$, and every prime $q\mid F(a,b)$ satisfies
\begin{equation}\label{ts-strong-cutoff}
 q\ge(\sqrt{2p}-1)^2>p,
 \qquad p\mid(q+1)^2-t_q(E_0)^2.
\end{equation}
\end{theorem}
\begin{proof}
We first specify the local inputs at $p$. For multiplicative reduction
over $\mathbb Q_p$, the full Tate representation is an extension
\[
 0\longrightarrow\mathbb F_p(\omega\eta)
 \longrightarrow E[p]\longrightarrow\mathbb F_p(\eta)
 \longrightarrow0,
\]
where $\omega$ is the mod-$p$ cyclotomic character and $\eta$ is the
unramified quadratic splitting character, possibly trivial.
For good supersingular reduction over $\mathbb Q_p$, the residual
representation is irreducible over $\mathbb F_p$. These local facts
are the Tate and good-reduction statements of
\cite[Propositions~2.12(b) and 2.11(c)]{DarmonDiamondTaylorLocal}.
For good ordinary reduction, the connected--etale filtration has
characters $\omega\lambda^{-1}$ and $\lambda$, where $\lambda$
is unramified. Its arithmetic Frobenius value modulo $p$ is the
unit root of $T^2-t_pT+p$, and hence is $t_p$ modulo $p$;
see \cite[\S1.3.1 and Exercise~1.30]{BoeckleLocalGalois}.
These are local statements over $\mathbb Q_p$, as in their proofs
by the Tate curve and finite-flat or formal groups over $\mathbb Z_p$.

Suppose $p\mid F$. The actual local calculation makes $p$ split in
$K$, identifies both completions with $\mathbb Q_p$, and gives
multiplicative reduction of $E_{a,b}$. Work at either completion.
The character $\nu$ is unramified and quadratic. From the Tate
subcharacter and \eqref{ts-module}, the boundary representation has
a subcharacter $\omega\eta\nu$ after coefficient extension. This
character has values in $\mathbb F_p$. Existence of an intertwiner
after extending coefficients implies existence over $\mathbb F_p$:
the intertwining maps are the kernel of linear equations over
$\mathbb F_p$, and scalar extension preserves this kernel. The
representations have finite image, so finitely many equations suffice.
Thus the supersingular irreducibility statement just cited rules out
supersingular reduction. No absolute-irreducibility upgrade of that
local statement has been assumed. The good boundary curve is ordinary.

Its Jordan--Holder characters after twisting by $\nu$ are
$\lambda\nu$ and $\omega\lambda^{-1}\nu$. The actual Tate
characters are $\eta$ and $\omega\eta$. Since $\omega$ is
nontrivial on inertia and all other characters here are unramified,
each pair has exactly one inertia-trivial character. The isomorphism
therefore gives
\begin{equation}\label{ts-unit-root}
 \lambda\nu=\eta,
 \qquad t_p(E_0)\equiv\pm1\pmod p.
\end{equation}
This compares the unique inertia-trivial Jordan--Holder character,
even if an extension splits. It does not take an unramified trace
of the full two-dimensional Tate representation at $p$.

By \eqref{ts-even-trace}, $t_p$ is even. But the Hasse bound and
$p>7$ give $|t_p|\le2\sqrt p<p-1$. The only integers in this
interval congruent to $1$ or $-1$ modulo $p$ are $1$ and $-1$,
which are odd. This contradiction proves $p\nmid F$.

Now every prime divisor $q$ of $F$ differs from $p$. The boundary comparison
of Theorem~\ref{bc-pure} gives $t_q\equiv\pm(q+1)\pmod p$.
Both $t_q$ and $q+1$ are even, so the corresponding nonzero
integer difference is divisible by $2p$. It is nonzero by the strict
Hasse inequality $|t_q|\le2\sqrt q<q+1$, and thus
\[
 2p\le |t_q\mp(q+1)|\le q+1+2\sqrt q=(\sqrt q+1)^2.
\]
This gives the first cutoff in \eqref{ts-strong-cutoff}. The strict
inequality with $p$ follows from $p+1>2\sqrt{2p}$ for $p>7$,
equivalently $p^2-6p+1>0$. The remaining divisibility is exactly
the boundary trace condition, now valid at every prime divisor.
\end{proof}

\begin{corollary}[Unconditional support of the actual pure branch]
\label{ts-pure-support}
Let $a,b$ be positive coprime integers, let $p>7$ be prime, and
suppose $F(a,b)=Q^p$ for an integer $Q>1$. Then $p\nmid Q$ and
every prime $q\mid Q$ satisfies \eqref{ts-strong-cutoff}. Neither $7$
nor $13$ divides $Q$, and
\begin{equation}\label{ts-root-size}
 Q\ge(\sqrt{2p}-1)^2>p.
\end{equation}
\end{corollary}
\begin{proof}
Theorem~\ref{bt-pure-module}, using the exact twist identity and its
stated modular dependencies, gives \eqref{ts-module} with a quadratic
character unramified outside $6$. Apply
Theorem~\ref{ts-boundary-support} and the previously proved boundary
exclusions at $7$ and $13$. Since $Q>1$ has a prime divisor, the
prime cutoff gives \eqref{ts-root-size}.
\end{proof}

For every sufficiently large fixed $p$, every prime factor now lies
in the permitted split rational-prime set of Theorem~\ref{bc-density}.
Its density is
\[
 \frac{2p^2-p-5}{2(p-1)^2(p+1)}.
\]
There is no exponent-prime exception. Removing the finite interval
below the cutoff does not change this fixed-$p$ density or supply
an error term uniform in $p$.

\begin{corollary}[Actual height lower bounds]\label{ts-height}
Under the hypotheses of Corollary~\ref{ts-pure-support}, put
$c=a+b$ and $H=\max(a,b)$. Then
\begin{align}
 \log c&\ge\frac p2\log(\sqrt{2p}-1),\label{ts-height-sum}\\
 \log H&\ge\frac p2\log(\sqrt{2p}-1)-\frac{\log13}{4}.
 \label{ts-height-max}
\end{align}
In particular $p^p<F(a,b)\le13H^4$.
\end{corollary}
\begin{proof}
The identity $4F=(a^2+b^2)^2+3(a+b)^4$ gives $F\le c^4$;
termwise comparison in the quartic gives $F\le13H^4$.
Combine these with $F=Q^p$ and \eqref{ts-root-size}, then take
logarithms. Its strict inequality $Q>p$ gives the last assertion.
\end{proof}

\paragraph{Verification and remaining scope.}
The full local module proof, the distinction between ordinary and
supersingular reduction, and the scalar-extension step have independent
ordinary reviews with the cited primary passages actually opened.
The elliptic representation inputs are not claimed to be formally
proved in Lean by arithmetic checks of the final inequalities.
These necessary height lower bounds remain compatible with arbitrarily
large seeds and moving large root primes. They do not give a
point-height upper bound, a contradiction to a cover's coefficient
budget, or an ABC proof. The finite local residue shadows remain
valid because they are not asserted to be global pure powers.
